\section{An exact finite spectral selection}
\label{sec:finite-selection}

The analytic argument in Lemma~\ref{lem:spectral-surplus} supplies all
spectral data needed in the proof of Theorem~\ref{thm:main}.  A supplementary
integer calculation gives a concrete selection in ambient dimension
\(16\), at the integer cutoff \(50000\).  This calculation checks the
finite spectral inequalities in~\eqref{eq:selection}; it does not construct
the finite control word or the metric numerically.  The analytic existence
argument remains independent of it.

The parameters of the calculation are
\begin{gather*}
 m=15,\qquad s=8,\qquad k=50000,\qquad q_*=\frac{101}{100},\\
 c=\frac{147}{1000},\qquad \alpha_*^2=\frac{99853}{100000},
 \qquad a_*^2=q_*^{1/15}\alpha_*^2.
\end{gather*}
The endpoint curvature inequality is strict because
\[
 1-\frac{2(q_*-1)}{m-1}-\alpha_*^2=\frac{29}{700000}>0.
\]
The integer inequality
\(101\cdot99853^{15}<100\cdot100000^{15}\)
shows that \(q_*(\alpha_*^2)^{15}<1\), and hence \(a_*<1\).
The decrease of the curvature threshold in the proof of
Lemma~\ref{lem:spectral-surplus} gives the margin for every
\(q\in[1,q_*]\).
The radial scale also meets its required smallness condition:
\[
 -\log a_*\le-\tfrac12\log\alpha_*^2
 <\frac{1-\alpha_*^2}{2\alpha_*^2}<\frac1{1000}.
\]
The cutoff integral in Proposition~\ref{prop:metric} is at least
\(\int_4^\infty(1+t)^{-5/4}\,dt=4\cdot5^{-1/4}>2\).
It therefore gives \(\mu<1/2000\), in particular \(b\ge-1/4\).
These are parameter checks; the control and transmission arguments are
still the analytic arguments of Sections~\ref{sec:control}--\ref{sec:conclusion}.

Here is the arithmetic bound underlying the selection.  Write
\(h=m-1\), \(Q=10^8\), and an angular eigenvalue as
\(\lambda=\mathrm{num}/\mathrm{den}>0\), with integer numerator and
denominator.  If \(\operatorname{isqrt}(a)=\lfloor\sqrt a\rfloor\) for a
nonnegative integer \(a\), define
\[
 S=\operatorname{isqrt}\!\left(
    \left\lfloor\frac{Q^2(h^2\mathrm{den}+4\mathrm{num})}
                         {\mathrm{den}}\right\rfloor\right)+1.
\]
Then \(S^2\mathrm{den}>Q^2(h^2\mathrm{den}+4\mathrm{num})\), so the
positive root of \(d(d+h)=\lambda\) satisfies the strict rational bound
\[
 d<\frac{S-hQ}{2Q}.
\]
For the charge indexed by \(b\), the exact eigenvalue data are
\[
 \mathrm{num}=100000\bigl(101l(l+14)-(2b-l)^2\bigr),
 \qquad \mathrm{den}=101\cdot99853.
\]
The multiplicities are the integers in~\eqref{eq:hopf-multiplicity}.
The two equal contributions indexed by \(b\) and \(l-b\) are combined,
except for the middle term when they coincide.  Increasing \(b\) from
\(0\) to \(\lfloor l/2\rfloor\) puts the eigenvalues, and therefore the
upper bounds for the exponents, in increasing order.

The supplied program \texttt{verification/counting-check.py}
uses these integer multiplicities and upper bounds.  It accepts an initial
group only when its upper-bounded exponent sum is strictly less than
\(k\) times its size.  It may take part of a repeated eigenvalue, which
corresponds to choosing that many vectors in its real eigenspace.  Low
degrees are selected in full whenever the exact inequality
\(l(l+14)<\alpha_*^2k(k+14)\) certifies all their roots to be below \(k\).
In odd degree the zero Hopf charge is absent, so the full-degree test
is a sufficient criterion rather than an exact characterization of the
largest root. The individual search ends when the smallest root is at least \(k\).
Thus every counted group has a true mean strictly below \(k\).

The resulting counts, using integer arithmetic, are
\begin{align*}
 \text{selected dimension}
   &=46785605044474340387088811257657817010714687366486217934111,\\
 h_{50000}(\mathbb R^{16})
   &=46779709349146362239266126538609505511855492134986686636251,\\
 \text{strict surplus}
   &=5895695327978147822684719048311498859195231499531297860.
\end{align*}
The selected dimension omits degrees \(0\) and \(1\), as required in
\eqref{eq:selection}.  Only the displayed strict integer comparison is
used by the certificate; a decimal ratio is unnecessary.

These parameters and selections satisfy all the finite-data hypotheses at
the end of Section~\ref{sec:counting}. Applying
Propositions~\ref{prop:metric} and~\ref{prop:transmission}, followed by the
growth argument in Section~\ref{sec:conclusion}, therefore gives the metric
of Theorem~\ref{thm:main} with \(n=16\) and \(k=50000\).
